\section{Scientific Graph \emojiscititle}
\subsection{Definition}
A molecule graph is a graphical representation of a molecule where atoms are represented as nodes and bonds between atoms are represented as edges.

\hy{Building a literature knowledge base towards transparent biomedical AI}

\subsection{RAG rational}
The underlying reason is that large language models (LLMs) lack domain-specific knowledge. Additionally, the molecules generated by the generative models are often not accurate enough. 

Molecule generation is crucial in fields like drug discovery and materials science, but traditional generative models often struggle with diversity and chemical validity. Retrieval-Augmented Generation (RAG)  improves diversity and accuracy by retrieving relevant molecules from a large database to guide the generation process. It ensures chemical validity by incorporating known valid structures, leverages existing knowledge for practical applications, and accelerates the generation process by narrowing the search space.

\subsection{Database Construction}
ChEMBL and ZINC


\subsection{Query/Task Instruction}
What kind of task or query should we encounter when GraphRAG on molecule graph.

Target-Specific Molecule Generation, molecular property prediction

\subsection{Retriever}
heuristic-based retriever
(To retrieve the nearest neighbor molecules from the database using cosine similarity.)

Learnable(
Retrieve ligands with high binding affinity.  The pre-trained PMINet
is reused here to extract interactive structural context information between the target protein P and ligands)

chain-of-thought 

knowledge graph
\subsection{Organizer}
Reranker(Scores Top K)
Information fusion. This module enables the retrieved exemplar molecules to modify the input
molecule towards the targeted design criteria. It achieves this by merging the embeddings of the
input and the retrieved exemplar molecules with a lightweight, trainable, standard cross attention
mechanism. (embedding fused and Diffusion-generation)

Incorporate this information into the molecule generation process through a cross-attention mechanism.

Inject retrieved information into the text prompt.
\subsection{Generator}
generative models

LLM, Diffusion model,transformer-based generative model


embedding-based fused

prompt-based Knowledge Injection
